%% kapa: exam=1 round=34 year=2023 period=2 form=A subject=accounting subject_ko=회계학 seq=2 number=42
%% source: Q-Net 감정평가사 기출문제 (cst00309) artlSeq=5211789 file=raw/1008/5211789/2023년도 제34회 감정평가사 1차시험 2교시.pdf
%% conversion: kordoc 4.24.51 -> tools/gen1.py (auto); review=proofread-sonnet; flags=none
\begin{question}{42}
무형자산에 관한 설명으로 옳지 않은 것은?
\begin{choices}
  \item 무형자산은 손상의 징후가 있거나 그 자산을 사용하지 않을 때에 상각을 중지한다.
  \item 무형자산의 인식기준을 충족하지 못해 비용으로 인식한 지출은 그 이후에 무형자산의 원가로 인식할 수 없다.
  \item 내부적으로 창출한 영업권은 자산으로 인식하지 아니한다.
  \item 개별취득 무형자산은 자산에서 발생하는 미래경제적효익의 유입가능성이 높다는 인식기준을 항상 충족한다.
  \item 무형자산으로 정의되려면 식별가능성, 자원에 대한 통제와 미래경제적효익의 존재를 충족하여야 한다.
\end{choices}
\end{question}
\answer{1}
